\label{sec:diferenciabilidad}

\textcolor{red}{definiciones derivabilidad}

\begin{definition}\textbf{Diferenciabilidad de campos escalares.} \label{def:dif_escalar}
\mbox{}
Sean $f: A \subset \Rn{n} \to \R$ y $\mathbf{x}_0 \in \interior{A}$. Decimos que $f$ es \emph{diferenciable en $\mathbf{x}_0$} si $\exists\:\mathbf{m} \in \Rn{n}$ tal que
\[
  \lim_{\mathbf{x} \to \mathbf{x}_0} \frac{f(\mathbf{x}) - f(\mathbf{x}_0) - \mathbf{m} \cdot (\mathbf{x} - \mathbf{x}_0)}{\norm{\mathbf{x} - \mathbf{x}_0}} = 0.
\]
Si $A$ es un conjunto abierto, decimos que $f$ es \emph{diferenciable} si es diferenciable en cada punto $\mathbf{x}_0 \in A$.
\end{definition}

\begin{definition}\textbf{Diferenciabilidad, caso general.} \label{def:dif_general}
\mbox{}

Sean $f: A \subset \Rn{n} \to \Rn{m}$ y $\mathbf{x}_0 \in \interior{(A)}$. Decimos que $f$ es \emph{diferenciable en $\mathbf{x}_0$} si $\exists M \in \Rnp{m}{n}$ tal que
\[
  \lim_{\mathbf{x} \to \mathbf{x}_0} \frac{\norm{f(\mathbf{x}) - f(\mathbf{x}_0) - M \cdot (\mathbf{x} - \mathbf{x}_0)}}{\norm{\mathbf{x} - \mathbf{x}_0}} = 0.
\]
Si $A$ es un conjunto abierto, decimos que $f$ es \emph{diferenciable} si es diferenciable en cada punto $\mathbf{x}_0 \in A$.
\end{definition}

\begin{theorem}\textbf{Diferenciabilidad y continuidad.} \label{teo:difCont}
\mbox{}

Sean $f: A \subset \Rn{n} \to \R$ y $\mathbf{x}_0 \in  A^{\circ}$. Si $f$ es diferenciable en $\mathbf{x}_0$, entonces $f$ es continua en $\mathbf{x}_0$.
 \begin{proof}
 \mbox{}
 
 Como $f$ es diferenciable en $\mathbf{x}_0, \exists\:\mathbf{m} \in \Rn{n}$ tal que
 \[
  \lim_{\mathbf{x} \to \mathbf{x}_0} \frac{f(\mathbf{x}) - f(\mathbf{x}_0) - \mathbf{m} \cdot (\mathbf{x} - \mathbf{x}_0)}{\norm{\mathbf{x} - \mathbf{x}_0}} = 0.
 \]
Por lo tanto, $\exists\:\delta_0 > 0$ tal que 
\[
 \mathbf{x} \in B_{\delta_0}^* \cap A \then 
 \abs{ \frac{f(\mathbf{x}) - f(\mathbf{x}_0) - \mathbf{m} \cdot (\mathbf{x} - \mathbf{x}_0)}{\norm{\mathbf{x} - \mathbf{x}_0}} } < 1.
\]
Entonces, para este $\delta_0$, se cumple que
\[
 \abs{ f(\mathbf{x}) - f(\mathbf{x}_0) - \mathbf{m} \cdot (\mathbf{x} - \mathbf{x}_0) } < \norm{\mathbf{x} - \mathbf{x}_0}, \, \forall \mathbf{x} \in B_{\delta_0}^* \cap A.
\]
% Como ambos miembros de esta desigualdad (estricta) son iguales cuando $\mathbf{x} = \mathbf{x}_0$, vale que
% \[
%  \abs{ f(\mathbf{x}) - f(\mathbf{x}_0) - \mathbf{m} \cdot (\mathbf{x} - \mathbf{x}_0) } \le \norm{\mathbf{x} - \mathbf{x}_0}, \, \forall \mathbf{x} \in B_{\delta_0} \cap A.
% \]
Por lo tanto,
\begin{align*}
\abs{ f(\mathbf{x}) - f(\mathbf{x}_0) } &=   \abs{ f(\mathbf{x}) - f(\mathbf{x}_0) - \mathbf{m} \cdot (\mathbf{x} - \mathbf{x}_0) + \mathbf{m} \cdot (\mathbf{x} - \mathbf{x}_0) } \le \\
                      &\le \abs{ f(\mathbf{x}) - f(\mathbf{x}_0) - \mathbf{m} \cdot (\mathbf{x} - \mathbf{x}_0) } + \abs{ \mathbf{m} \cdot (\mathbf{x} - \mathbf{x}_0) } < \\
                      &< \norm{\mathbf{x} - \mathbf{x}_0} + \abs{ \mathbf{m} \cdot (\mathbf{x} - \mathbf{x}_0) }, \, 
                      \forall \mathbf{x} \in B^*_{\delta_0} \cap A.
\end{align*}
Por la desigualdad de Cauchy-Schwarz, $\abs{ \mathbf{m} \cdot (\mathbf{x} - \mathbf{x}_0) } \le \norm{\mathbf{m}} \norm{\mathbf{x} - \mathbf{x}_0}$, y entonces:
\begin{align*}
\abs{ f(\mathbf{x}) - f(\mathbf{x}_0) } &< \norm{\mathbf{x} - \mathbf{x}_0} + \abs{ \mathbf{m} \cdot (\mathbf{x} - \mathbf{x}_0) } \le \\
                      &\le \norm{\mathbf{x} - \mathbf{x}_0} + \norm{\mathbf{m}} \norm{\mathbf{x} - \mathbf{x}_0} = \\
                      &=  \left( 1 + \norm{\mathbf{m}} \right) \norm{\mathbf{x} - \mathbf{x}_0}, \, 
                      \forall \mathbf{x} \in B^*_{\delta_0} \cap A.
\end{align*}
Esta condici\'on garantiza la continuidad de $f$ en $\mathbf{x}_0$. En efecto, dado $\epsilon > 0$, sea $\delta = \min \{\delta_0, \frac{\epsilon}{1 + \norm{\mathbf{m}}} \}$. Para este $\delta$ se cumple que:
\begin{align*}
 \abs{ f(\mathbf{x}) - f(\mathbf{x}_0) } &< \left( 1 + \norm{\mathbf{m}} \right) \norm{\mathbf{x} - \mathbf{x}_0} < \\
                       &<   \left( 1 + \norm{\mathbf{m}} \right) \delta \le \\
                       &\le \left( 1 + \norm{\mathbf{m}} \right) \frac{\epsilon}{1 + \norm{\mathbf{m}}} = \epsilon, 
                       \, \forall \mathbf{x} \in B^*_{\delta} \cap A,
\end{align*}
lo que, por definici\'on de l\'imite \footnote{Como $\mathbf{x}_0 \in \interior{(A)}$, $\mathbf{x}_0$ es un punto de acumulaci\'on de $A$. }, significa que 
\[
 \lim_{\mathbf{x} \to \mathbf{x}_0}f(\mathbf{x}) = f(\mathbf{x}_0).
\]
Es decir, $f$ es continua en $\mathbf{x}_0$.
\end{proof}
\begin{obs} La implicaci\'on rec\'iproca de la proposici\'on \eqref{teo:difCont} es falsa. 
 \begin{proof} La funci\'on
  \begin{align*}
      g & :\R \to \R \\
        & g(t) = \abs{t}
 \end{align*}
es continua en $t_0 = 0$, pero no es diferenciable en $t_0 = 0$. Se deja como ejercicio para el lector verificar esta afirmaci\'on.
 \end{proof}
\end{obs}
\end{theorem}

\begin{theorem}\textbf{Diferenciabilidad y existencia de derivadas parciales.} \label{teo:difDer}
\mbox{}

\iffalse
Sean $f: A \subset \Rn{2} \to \R$ y $(\mathbf{x}_0, y_0) \in \interior{(A)}$, tales que $\exists \alpha, \beta \in \R \, /$
\[
 \lim_{(\mathbf{x},y) \to (\mathbf{x}_0,y_0)} \frac{f(\mathbf{x},y) - f(\mathbf{x}_0,y_0) - \alpha(\mathbf{x} - \mathbf{x}_0) - \beta(y - y_0)}
                                 {\norm{(\mathbf{x} - \mathbf{x}_0, y - y_0)}} = 0
\]
(es decir, $f$ es diferenciable en $(\mathbf{x}_0,y_0)$). Entonces existen las derivadas parciales de $f$ en $(\mathbf{x}_0,y_0)$ y, adem\'as,
\[
 \dif{f}{\mathbf{x}} (\mathbf{x}_0,y_0) = \alpha, \quad \dif{f}{y}(\mathbf{x}_0,y_0) = \beta.
\]

 \begin{proof}
 \mbox{}
 
Sea $(\mathbf{x},y) = (\mathbf{x}_0,y_0) + h \mathbf{e}_1$, siendo $h \in \R$ y $\mathbf{e}_1$ el primer vector can\'onico ($\mathbf{e}_1 = (1,0)$). Es decir, definimos una funci\'on $\mathbf{g}:(-\delta,\delta) \to \Rn{2} / \mathbf{g}(h) = (\mathbf{x}_0,y_0) + h \mathbf{e}_1$, con $\delta > 0$, lo suficientemente peque\~{n}o como para que $\text{Img}(\mathbf{g}) \subset A$ \footnote{Cu\'al de las hip\'oteis asegura la existencia de un tal $\delta$?}, y llamamos $(\mathbf{x},y) = \mathbf{g}(h)$. Notemos que $\mathbf{g}(h) \to (\mathbf{x}_0,y_0)$ cuando $h \to 0$. 

Por la diferenciabilidad, tenemos que 
\[
 \lim_{(\mathbf{x},y) \to (\mathbf{x}_0,y_0)} \frac{f(\mathbf{x},y) - f(\mathbf{x}_0,y_0) - \alpha(\mathbf{x} - \mathbf{x}_0) - \beta(y - y_0)}
                                 {\norm{(\mathbf{x} - \mathbf{x}_0, y - y_0)}} = 0
\]
y entonces, si componemos con la funci\'on $\mathbf{g}$ (reemplazamos $(\mathbf{x},y) = (\mathbf{x}_0,y_0) + h \mathbf{e}_1$), la propiedad del l\'imite de la composici\'on \eqref{prop:lim_comp} implica que
\begin{align*}
  0 & = \lim_{h \to 0} \frac{f(\mathbf{x}_0 + h,y_0) - f(\mathbf{x}_0,y_0) - \alpha(\mathbf{x}_0 + h - \mathbf{x}_0) - \beta(y_0 - y_0)}
                                 {\norm{(\mathbf{x}_0 + h - \mathbf{x}_0, y_0 - y_0)}} = \\
    & = \lim_{h \to 0} \frac{f(\mathbf{x}_0 + h,y_0) - f(\mathbf{x}_0,y_0) - \alpha h}{\norm{h \mathbf{e}_1}} = \\
    & = \lim_{h \to 0} \frac{f(\mathbf{x}_0 + h,y_0) - f(\mathbf{x}_0,y_0) - \alpha h}{\abs{h}} .
\end{align*}                                
Ahora bien, este l\'imite es $0$ si y s\'olo si
\[
 \lim_{h \to 0} \abs{ \frac{f(\mathbf{x}_0 + h,y_0) - f(\mathbf{x}_0,y_0) - \alpha h}{\abs{h}} } = 0,
\]
por la propiedad \eqref{prop:lim_1}. Entonces,
\begin{align*}
  0 & =  \lim_{h \to 0} \abs{ \frac{f(\mathbf{x}_0 + h,y_0) - f(\mathbf{x}_0,y_0) - \alpha h}{\abs{h}} }
      = \lim_{h \to 0} \frac{ \abs{ f(\mathbf{x}_0 + h,y_0) - f(\mathbf{x}_0,y_0) - \alpha h }}{\abs{ \, \abs{h} \, }} = \\
    & = \lim_{h \to 0} \frac{ \abs{ f(\mathbf{x}_0 + h,y_0) - f(\mathbf{x}_0,y_0) - \alpha h }}{\abs{h}}
      = \lim_{h \to 0} \abs{ \frac{ f(\mathbf{x}_0 + h,y_0) - f(\mathbf{x}_0,y_0) - \alpha h }{h}},
\end{align*} 
lo cual, otra vez por la propiedad \eqref{prop:lim_1}, implica que 
\[
 0 = \lim_{h \to 0} \frac{ f(\mathbf{x}_0 + h,y_0) - f(\mathbf{x}_0,y_0) - \alpha h }{h} =
     \lim_{h \to 0} \left( \frac{ f(\mathbf{x}_0 + h,y_0) - f(\mathbf{x}_0,y_0) }{h} - \alpha \right) .
\]
\fi

\iffalse
Entonces, por la propiedad \eqref{prop:lim_2}, tenemos que
\[
 \lim_{h \to 0} \frac{ f(\mathbf{x}_0 + h,y_0) - f(\mathbf{x}_0,y_0) }{h} = \alpha,
\]
\fi
Sean $f: A \subset \Rn{n} \to \R$ y $\mathbf{x}_0 \in \interior{(A)}$, tales que $\exists \mathbf{m} = (m_1, m_2, \cdots, m_n) \in \Rn{n} \, /$
\[
 \lim_{\mathbf{x} \to \mathbf{x}_0} \frac{f(\mathbf{x}) - f(\mathbf{x}_0) - \mathbf{m} \cdot (\mathbf{x} - \mathbf{x}_0)}
                                 {\norm{\mathbf{x} - \mathbf{x}_0}} = 0
\]
(es decir, $f$ es diferenciable en $\mathbf{x}_0$). Entonces existen las derivadas parciales de $f$ en $\mathbf{x}_0$ y, adem\'as,
\[
 \dif{f}{x_j} (\mathbf{x}_0) = m_j \quad \forall j = 1, 2, \cdots, n.
\]

 \begin{proof}
 \mbox{}
 
Sea $\mathbf{x} = \mathbf{x}_0 + h \mathbf{e}_j$, siendo $h \in \R$ y $\mathbf{e}_j$ el $j$-\'esimo vector can\'onico (el vector de $\Rn{n}$ que tiene un $1$ en la posici\'on $j$ y ceros en todas las dem\'as). Es decir, definimos una funci\'on $\mathbf{g}:(-\delta,\delta) \to \Rn{n}\;|\; \mathbf{g}(h) = \mathbf{x}_0 + h \mathbf{e}_j$, con $\delta > 0$, lo suficientemente peque\~{n}o como para que $\text{Img}(\mathbf{g}) \subset A$ \footnote{¿Cul\'al de las hip\'otesis asegura la existencia de un tal $\delta$?}, y llamamos $\mathbf{x} = \mathbf{g}(h)$. Notemos que $\mathbf{g}(h) \to \mathbf{x}_0$ cuando $h \to 0$. 

Por la diferenciabilidad, tenemos que 
\[
 \lim_{\mathbf{x} \to \mathbf{x}_0} \frac{f(\mathbf{x}) - f(\mathbf{x}_0) - \mathbf{m} \cdot (\mathbf{x} - \mathbf{x}_0)}
                                 {\norm{\mathbf{x} - \mathbf{x}_0}} = 0
\]
y entonces, si componemos con la funci\'on $\mathbf{g}$ (reemplazamos $\mathbf{x} = \mathbf{x}_0 + h \mathbf{e}_j$), la propiedad del l\'imite de la composici\'on \eqref{prop:lim_comp} implica que
\begin{align*}
  0 & = \lim_{h \to 0} \frac{f(\mathbf{x}_0 + h \mathbf{e}_j) - f(\mathbf{x}_0) - \mathbf{m} \cdot ((\mathbf{x}_0 + h\mathbf{e}_j) - \mathbf{x}_0)}
                                 {\norm{(\mathbf{x}_0 + h\mathbf{e}_j) - \mathbf{x}_0}} = \\
    & = \lim_{h \to 0} \frac{f(\mathbf{x}_0 + h \mathbf{e}_j) - f(\mathbf{x}_0) - \mathbf{m} \cdot (h\mathbf{e}_j)}{\norm{h \mathbf{e}_j}} = \lim_{h \to 0} \frac{f(\mathbf{x}_0 + h \mathbf{e}_j) - f(\mathbf{x}_0) - h(\mathbf{m} \cdot \mathbf{e}_j)}{\norm{h \mathbf{e}_j}} =\\
    & = \lim_{h \to 0} \frac{f(\mathbf{x}_0 + h \mathbf{e}_j) - f(\mathbf{x}_0) - h m_j}{\abs{h}} .
\end{align*}                                
Ahora bien, este l\'imite es $0$ si y s\'olo si
\[
 \lim_{h \to 0} \abs{ \frac{f(\mathbf{x}_0 + h \mathbf{e}_j) - f(\mathbf{x}_0) - h m_j}{\abs{h}} } = 0,
\]
por la propiedad \eqref{prop:lim_1}. Entonces,
\begin{align*}
  0 & = \lim_{h \to 0} \abs{ \frac{f(\mathbf{x}_0 + h \mathbf{e}_j) - f(\mathbf{x}_0) - h m_j}{\abs{h}} }
      = \lim_{h \to 0} \frac{\abs{f(\mathbf{x}_0 + h \mathbf{e}_j) - f(\mathbf{x}_0) - h m_j}}{\abs{\abs{h}}} = \\
    & = \lim_{h \to 0} \frac{\abs{f(\mathbf{x}_0 + h \mathbf{e}_j) - f(\mathbf{x}_0) - h m_j}}{\abs{h}}
      = \lim_{h \to 0} \abs{ \frac{f(\mathbf{x}_0 + h \mathbf{e}_j) - f(\mathbf{x}_0) - h m_j}{h} },
\end{align*} 
lo cual, otra vez por la propiedad \eqref{prop:lim_1}, implica que 
\[
 0 = \lim_{h \to 0} \frac{f(\mathbf{x}_0 + h \mathbf{e}_j) - f(\mathbf{x}_0) - h m_j}{h} =
     \lim_{h \to 0} \left( \frac{f(\mathbf{x}_0 + h \mathbf{e}_j) - f(\mathbf{x}_0)}{h} - m_j \right) .
\]

Es f\'acil probar por definici\'on que esta \'ultima condici\'on implica que
\[
 \lim_{h \to 0} \frac{ f(\mathbf{x}_0 + h \mathbf{e}_j) - f(\mathbf{x}_0) }{h} = m_j.
\]
En efecto, dado $\epsilon > 0$ existe $\delta > 0$ tal que 
\[
 \abs{ \left( \frac{ f(\mathbf{x}_0 + h \mathbf{e}_j) - f(\mathbf{x}_0) }{h} - m_j \right) - 0 } < \epsilon, \, \forall h : 0 < \abs{h} < \delta,
\]
porque el l\'imite es 0. Como claramente 
\[
 \abs{ \left( \frac{ f(\mathbf{x}_0 + h \mathbf{e}_j) - f(\mathbf{x}_0) }{h} - m_j \right) - 0 } = 
 \abs{\frac{ f(\mathbf{x}_0 + h \mathbf{e}_j) - f(\mathbf{x}_0) }{h} - m_j},
\]
se tiene (por definici\'on de l\'imite) que 
\[
 \lim_{h \to 0} \frac{ f(\mathbf{x}_0 + h \mathbf{e}_j) - f(\mathbf{x}_0) }{h} = m_j.
\]
Esto \'ultimo, por definici\'on de derivada parcial, significa que
\[
  \dif{f}{x_j} (\mathbf{x}_0) = m_j,
\]
como quer\'iamos probar. 

 \end{proof}

\begin{obs}
 Es un error decir, en el \'ultimo paso de la demostraci\'on, que por la linealidad del l\'imite
 \[
  \lim_{h \to 0} \left( \frac{f(\mathbf{x}_0 + h \mathbf{e}_j) - f(\mathbf{x}_0)}{h} - m_j \right) = 
  \lim_{h \to 0} \left( \frac{f(\mathbf{x}_0 + h \mathbf{e}_j) - f(\mathbf{x}_0)}{h} \right) - \lim_{h \to 0} m_j = 0,
 \]
ya que la existencia del l\'imite 
\[
\lim_{h \to 0} \left( \frac{f(\mathbf{x}_0 + h \mathbf{e}_j) - f(\mathbf{x}_0)}{h} \right)
\]
es parte de lo que debemos probar.
\end{obs}
\begin{obs} La implicaci\'on rec\'iproca de la proposici\'on \eqref{teo:difDer} es falsa.
\begin{proof}
\mbox{}
  
La funci\'on
 \[
       f :\Rn{2} \to \R
      \]
      \[
        f(x,y) = 
        \begin{cases}
         0 & \text{ si } x y = 0  \\
         1 & \text{ en otro caso }
        \end{cases}
\]
tiene derivadas parciales en el origen, pero no es diferenciable en el origen.\\
Se deja como ejercicio para el lector la demostraci\'on de esta afirmaci\'on.
\end{proof}
\end{obs}
\end{theorem}

Una consecuencia directa del teorema \eqref{teo:difDer} es el siguiente corolario.
\begin{corollary}
 Sean $f: A \subset \Rn{n} \to \R$ y $\mathbf{x}_0 \in A^{\circ}$. Son equivalentes:
 \begin{enumerate} %[I.]
  \item $f$ es diferenciable en $\mathbf{x}_0$;
  \item existen todas las derivadas parciales de $f$ en $\mathbf{x}_0$ y 
  \[
   \lim_{\mathbf{x} \to \mathbf{x}_0} \frac{f(\mathbf{x}) - f(\mathbf{x}_0) - \grad f (\mathbf{x}_0) \cdot (\mathbf{x} - \mathbf{x}_0)}{\norm{\mathbf{x} - \mathbf{x}_0}} = 0.
  \]
 \end{enumerate}
\end{corollary}

\begin{theorem}\textbf{Diferenciabilidad y existencia de derivadas direccionales.}\label{teo:difDir}
% La idea es incluir tres demostraciones, una usando la misma idea que para la demo de las derivadas parciales, otra usando que el delta f se puede escribir de una determinada manera y construyendo el cociente incremental, una tercera utilizando la regla de la cadena
\mbox{}

Sea $f: A \subset \Rn{n} \to \R$ diferenciable en $\mathbf{x}_0 \in \interior{A}$. Entonces, $\forall \hat{\mathbf{u}} \in \Rn{n}\;\text{tal que } \norm{\hat{\mathbf{u}}} = 1$ existe en $\mathbf{x}_0$ la derivada de $f$ en la direcci\'on de $\hat{\mathbf{u}}$ y, adem\'as,
\[
 \frac{\partial f}{\partial \hat{\mathbf{u}}} (\mathbf{x}_0) = \grad f (\mathbf{x}_0) \cdot \hat{\mathbf{u}}.
\]

 \begin{proof}
 \mbox{}
 
 Sea $\mathbf{x} = \mathbf{x}_0 + h \hat{\mathbf{u}}$, siendo $h \in \R$. Es decir, definimos una funci\'on 
 \begin{gather*}
      \mathbf{g} :(-\delta,\delta) \to \Rn{n} \\
      \mathbf{g}(h) = \mathbf{x}_0 + h \hat{\mathbf{u}}
 \end{gather*}

con $\delta > 0$, lo suficientemente peque\~{n}o como para que $\text{Img}(\mathbf{g}) \subset A$, y llamamos $\mathbf{x} = \mathbf{g}(h)$. Notemos que $\mathbf{g}(h) \to \mathbf{x}_0$ cuando $h \to 0$. 

Por la diferenciabilidad, tenemos que 
\[
 \lim_{\mathbf{x} \to \mathbf{x}_0} \frac{f(\mathbf{x}) - f(\mathbf{x}_0) - \grad f (\mathbf{x}_0) \cdot (\mathbf{x} - \mathbf{x}_0)}
                                 {\norm{\mathbf{x} - \mathbf{x}_0}} = 0
\]
y entonces, si componemos con la funci\'on $\mathbf{g}$ (reemplazamos $\mathbf{x} = \mathbf{x}_0 + h \hat{\mathbf{u}}$), la propiedad del l\'imite de la composici\'on \eqref{prop:lim_comp} implica que
\begin{align*}
  0 & = \lim_{h \to 0} \frac{f(\mathbf{x}_0 + h \hat{\mathbf{u}}) - f(\mathbf{x}_0) - \grad f (\mathbf{x}_0) \cdot ((\mathbf{x}_0 + h _{\delta}(\mathbf{x}_0)\hat{\mathbf{u}}) - \mathbf{x}_0) }
                    {\norm{(\mathbf{x}_0 + h \hat{\mathbf{u}}) - \mathbf{x}_0}} = \\
    & = \lim_{h \to 0} \frac{f(\mathbf{x}_0 + h \hat{\mathbf{u}}) - f(\mathbf{x}_0) - \grad f (\mathbf{x}_0) \cdot (h \hat{\mathbf{u}}) }
                    {\norm{h \hat{\mathbf{u}}}} = \\
    & = \lim_{h \to 0} \frac{f(\mathbf{x}_0 + h \hat{\mathbf{u}}) - f(\mathbf{x}_0) - h \grad f (\mathbf{x}_0) \cdot \hat{\mathbf{u}} }
                    {\abs{h} \norm{\hat{\mathbf{u}}}} = \\
    & = \lim_{h \to 0} \frac{f(\mathbf{x}_0 + h \hat{\mathbf{u}}) - f(\mathbf{x}_0) - h \grad f (\mathbf{x}_0) \cdot \hat{\mathbf{u}} }
                    {\abs{h}}.  
\end{align*}                                
Como en la demostraci\'on de la propiedad \eqref{teo:difDer}, este l\'imite es $0$ si y s\'olo si
\[
 0 = \lim_{h \to 0} \frac{f(\mathbf{x}_0 + h \hat{\mathbf{u}}) - f(\mathbf{x}_0) - h \grad f (\mathbf{x}_0) \cdot \hat{\mathbf{u}}}{h} = 
     \lim_{h \to 0} \left( \frac{f(\mathbf{x}_0 + h \hat{\mathbf{u}}) - f(\mathbf{x}_0)}{h} - \grad f (\mathbf{x}_0) \cdot \hat{\mathbf{u}} \right),
\]
que, como en la demostraci\'on de \eqref{teo:difDer}, implica que 
\[
 \lim_{h \to 0} \frac{f(\mathbf{x}_0 + h \hat{\mathbf{u}}) - f(\mathbf{x}_0)}{h} = \grad f (\mathbf{x}_0) \cdot \hat{\mathbf{u}}
\]
y por lo tanto, por definici\'on,
\[
 \frac{\partial f}{\partial \hat{\mathbf{u}}} (\mathbf{x}_0) = \grad f (\mathbf{x}_0) \cdot \hat{\mathbf{u}}.
\]
 \end{proof}
\end{theorem}

\begin{theorem}\textbf{Valores extremos de la derivada direccional.} \label{teo:max_deriv}
\mbox{}

 Sean $f: A \subset \Rn{n} \to \R$ diferenciable en un punto $\mathbf{x}_0 \in \interior{A}$. Entonces 
 \[
  \max_{\norm{\hat{\mathbf{u}}} = 1} \frac{\partial f}{\partial \hat{\mathbf{u}}} (\mathbf{x}_0) = \norm{\grad f (\mathbf{x}_0)}, \quad 
  \min_{\norm{\hat{\mathbf{u}}} = 1} \frac{\partial f}{\partial \hat{\mathbf{u}}} (\mathbf{x}_0) = -\norm{\grad f (\mathbf{x}_0)}
 \]
y, si $\grad f(\mathbf{x}_0) \ne \mathbf{0}$,
\[
 \frac{\partial f}{\partial \hat{\mathbf{u}}} (\mathbf{x}_0) = \norm{\grad f (\mathbf{x}_0)} \iff \hat{\mathbf{u}} = \frac{\grad f(\mathbf{x}_0)}{\norm{\grad f(\mathbf{x}_0)}}
\]
y
\[
 \frac{\partial f}{\partial \hat{\mathbf{u}}} (\mathbf{x}_0) = -\norm{\grad f (\mathbf{x}_0)} \iff \hat{\mathbf{u}} = -\frac{\grad f(\mathbf{x}_0)}{\norm{\grad f(\mathbf{x}_0)}}.
\]
\begin{proof}
\mbox{}

 Como $f$ es diferenciable en $\mathbf{x}_0$, por el teorema \eqref{teo:difDir}, $\forall \hat{\mathbf{u}} \in \Rn{n}\;\text{tal que } \norm{\hat{\mathbf{u}}} = 1$ existe en $\mathbf{x}_0$ la derivada de $f$ en la direcci\'on de $\hat{\mathbf{u}}$ y, adem\'as,
\[
 \frac{\partial f}{\partial \hat{\mathbf{u}}} (\mathbf{x}_0) = \grad f (\mathbf{x}_0) \cdot \hat{\mathbf{u}}.
\]
Por la desigualdad de Cauchy-Schwarz, 
\[
 \abs{\frac{\partial f}{\partial \hat{\mathbf{u}}} (\mathbf{x}_0)} = \abs{\grad f (\mathbf{x}_0) \cdot \hat{\mathbf{u}}} \le \norm{\grad f (\mathbf{x}_0)}\norm{\hat{\mathbf{u}}} = \norm{\grad f (\mathbf{x}_0)},
\]
de modo que tenemos las cotas:
\[
 -\norm{\grad f (\mathbf{x}_0)} \le \frac{\partial f}{\partial \hat{\mathbf{u}}}(\mathbf{x}_0) \le \norm{\grad f (\mathbf{x}_0)}.
\]
Adem\'as, la desigualdad de Cauchy-Schwarz nos dice que vale
\[
 \abs{\frac{\partial f}{\partial \hat{\mathbf{u}}} (\mathbf{x}_0)} = \norm{\grad f (\mathbf{x}_0)}
\]
si y s\'olo si el conjunto $\{\grad f(\mathbf{x}_0), \hat{\mathbf{u}} \}$ es linealmente dependiente, condici\'on que, si $\grad f(\mathbf{x}_0) \ne \mathbf{0}$, es equivalente a decir que $\hat{\mathbf{u}}$ es un vector unitario en la direcci\'on de $\grad f(\mathbf{x}_0)$. Existen exactamente dos vectores $\hat{\mathbf{u}}$ con esta caracter\'istica:
\[
 \hat{\mathbf{u}}_1 = \frac{\grad f(\mathbf{x}_0)}{\norm{\grad f(\mathbf{x}_0)}} \quad \text{ y } \quad \hat{\mathbf{u}}_2 = -\frac{\grad f(\mathbf{x}_0)}{\norm{\grad f(\mathbf{x}_0)}}.
\]
Por c\'alculo directo:
\begin{align*}
 \frac{\partial f}{\partial \hat{\mathbf{u}}_1} (\mathbf{x}_0) &= \grad f (\mathbf{x}_0) \cdot \frac{\grad f(\mathbf{x}_0)}{\norm{\grad f(\mathbf{x}_0)}} = \frac{ \grad f (\mathbf{x}_0) \cdot \grad f(\mathbf{x}_0) }{\norm{\grad f(\mathbf{x}_0)}} = \frac{\norm{\grad f (\mathbf{x}_0)}^2 }{\norm{\grad f(\mathbf{x}_0)}} = \norm{\grad f (\mathbf{x}_0)} \\
%  
 \frac{\partial f}{\partial \hat{\mathbf{u}}_2} (\mathbf{x}_0) &= \grad f (\mathbf{x}_0) \cdot \left( - \frac{\grad f(\mathbf{x}_0)}{\norm{\grad f(\mathbf{x}_0)}} \right) = -\frac{ \grad f (\mathbf{x}_0) \cdot \grad f(\mathbf{x}_0) }{\norm{\grad f(\mathbf{x}_0)}} = -\frac{\norm{\grad f (\mathbf{x}_0)}^2 }{\norm{\grad f(\mathbf{x}_0)}} = -\norm{\grad f (\mathbf{x}_0)},
\end{align*}
de modo que efectivamente
\[
 \max_{\norm{\hat{\mathbf{u}}}=1} \frac{\partial f}{\partial \hat{\mathbf{u}}} (\mathbf{x}_0) = \norm{\grad f (\mathbf{x}_0)}
\]
y el m\'aximo se realiza en $\hat{\mathbf{u}}_1 = \frac{\grad f(\mathbf{x}_0)}{\norm{\grad f(\mathbf{x}_0)}}$,
y
\[
 \min_{\norm{\hat{\mathbf{u}}}=1} \frac{\partial f}{\partial \hat{\mathbf{u}}} (\mathbf{x}_0) = -\norm{\grad f (\mathbf{x}_0)}
\]
y el m\'inimo se realiza en $\hat{\mathbf{u}}_2 = -\frac{\grad f(\mathbf{x}_0)}{\norm{\grad f(\mathbf{x}_0)}}$. \\

Finalmente, observemos que, si $\grad f(\mathbf{x}_0) = \mathbf{0}$, 
\[
 \frac{\partial f}{\partial \hat{\mathbf{u}}} (\mathbf{x}_0) = \grad f (\mathbf{x}_0) \cdot \hat{\mathbf{u}} = 0 = \norm{\grad f(\mathbf{x}_0)} \quad \forall \hat{\mathbf{u}} \in \Rn{n},
\]
y por lo tanto:
\[
 \max_{\norm{\hat{\mathbf{u}}}=1} \frac{\partial f}{\partial \hat{\mathbf{u}}} (\mathbf{x}_0) = \norm{\grad f(\mathbf{x}_0)} = 0
\]
y
\[
 \min_{\norm{\hat{\mathbf{u}}}=1} \frac{\partial f}{\partial \hat{\mathbf{u}}} (\mathbf{x}_0) = -\norm{\grad f(\mathbf{x}_0)} = 0.
\]
\end{proof}
\end{theorem}

\begin{definition}
  Sea $\mathbf{F}:U\subseteq\Rn{n}\to\Rn{m}$, diferenciable en $\mathbf{x}_0\in\interior{U}$. Se define la \emph{matriz diferencial de $\mathbf{F}$ en $\mathbf{x}_0$}, y se nota por $\boldsymbol{D}_{\mathbf{F}}(\mathbf{x}_0)$, a
  \textcolor{red}{agregar X0} 
  \[
  \boldsymbol{D}_{\mathbf{F}}(\mathbf{x}_0)=
    \begin{bmatrix} 
      \grad{f_1} \\
      \grad{f_2} \\
      \vdots     \\
      \grad{f_m}
    \end{bmatrix}
    =
    \begin{bmatrix} 
      \dif{f_1}{x_1} & \dif{f_1}{x_2} & \cdots & \dif{f_1}{x_n} \\
                     &                &        &                \\
      \dif{f_2}{x_1} & \dif{f_2}{x_2} & \cdots & \dif{f_2}{x_n} \\
                     &                &        &                \\
       \vdots         &  \vdots        & \ddots & \vdots         \\
                     &                &        &                \\
      \dif{f_m}{x_1} & \dif{f_m}{x_2} & \cdots & \dif{f_m}{x_n}
    \end{bmatrix},
  \]
  donde $\mathbf{F}=(f_1, f_2, \dots, f_m)\;|\;f_i:U\subseteq\Rn{n}\to\R,\;i=1,2,\dots, m.$

  \textcolor{red}{hacer otra definicion}
  Y a su vez, en el caso particular en que $n=m$; se define el \emph{jacobiano de $\mathbf{F}$ en $\mathbf{x}_0$}, y se lo nota por $J_{\mathbf{F}}(\mathbf{x}_0)$, al determinante de la matriz diferencial, \emph{i.e.}
  \[
    J_{\mathbf{F}}(\mathbf{x}_0) = \det(\boldsymbol{D}_{\mathbf{F}}(\mathbf{x}_0)). 
  \]
\end{definition}

\begin{obs}
Notar que $\boldsymbol{D}_{\mathbf{F}}\in\R^{m\times n}$.
\end{obs}

\begin{theorem}\textbf{Regla de la cadena.} \label{teo:cadena}
\mbox{}
\textcolor{red}{campos en mayus}
  Sean $\mathbf{f}, \mathbf{g}$ funciones 
    \begin{align*}
    \mathbf{g} &:A \subset \Rn{n} \to B \subset \Rn{m} \\
    \mathbf{f} &:B \subset \Rn{m} \to \Rn{p},
    \end{align*}
  y un punto $\mathbf{x}_0 \in \interior{A} \text{ tal que } \mathbf{g}(\mathbf{x}_0) \in \interior{B}$. Supongamos que $\mathbf{g}$ es diferenciable en $\mathbf{x}_0$ y que $\mathbf{f}$ es diferenciable en $\mathbf{y}_0 = \mathbf{g}(\mathbf{x}_0)$. Entonces la composic\'on $\mathbf{f} \circ \mathbf{g}$ es diferenciable en $\mathbf{x}_0$ y, adem\'as,
  \[
   \boldsymbol{D}_{\mathbf{f} \circ \mathbf{g}}{(\mathbf{x}_0)} = \boldsymbol{D}_\mathbf{f}{(\mathbf{y}_0)} 
   \boldsymbol{D}_\mathbf{g}{(\mathbf{x}_0)}.
  \]
\end{theorem}

\begin{theorem}\textbf{Teorema de la funci\'on inversa.} \label{teo:inversa}
\mbox{}

 Sean $A \subset \Rn{n}$ un conjunto abierto y $\mathbf{F}: A \subset \Rn{n} \to \Rn{n}$ una funci\'on de clase $\mathcal{C}^1$. Sea $\mathbf{x}_0 \in A$ y supongamos que $J_\mathbf{F}(\mathbf{x}) \ne 0$. Entonces existen un entorno abierto $U \subset A$ de $\mathbf{x}_0$ y un entorno abierto $V \subset \Rn{n}$ de $\mathbf{F}(\mathbf{x}_0)$ tal que $\mathbf{F}:U \to V$ (la restricci\'on de $\mathbf{F}$ a $U$) tiene inversa $\mathbf{F}^{-1}:V \to U$. Esta funci\'on $\mathbf{F}^{-1}$ es de clase $\mathcal{C}^1$ y, adem\'as,
 \[
  \boldsymbol{D}_{\mathbf{F}^{-1}}{(\mathbf{F}(\mathbf{x}_0))} = \boldsymbol{D}_\mathbf{F}{(\mathbf{x}_0)}^{-1}.
 \]
 Si $\mathbf{F}$ es de clase $\mathcal{C}^p, p \ge 1,$ entonces $\mathbf{F}^{-1}$ tambi\'en.
\end{theorem}

\begin{theorem}\textbf{Caso particular del teorema de la funci\'on impl\'icita.} \label{teo:implicit_part}
\mbox{}

Sean $A \subset \Rn{3}$ un conjunto abierto y $F: A \to \R$ una funci\'on de clase $\mathcal{C}^1$. Sea $(x_0,y_0,z_0) \in A$ tal que $F(x_0,y_0,z_0) = 0$. Supongamos que 
 \[
    \dif{F}{z} (x_0,y_0,z_0) \ne 0.
 \]
 Entonces existen un entorno abierto $U \subset \Rn{2}$ de $(x_0,y_0)$
 %, un entorno abierto $V \subset \Rn{m}$ de $y_0$ 
 y una \'unica funci\'on $g: U \to \R$ tales que $g(x_0,y_0) = z_0$ y
 \[
  F(x,y,g(x,y)) = 0,\quad \forall (x,y) \in U.
 \]
 M\'as a\'un, $g$ es de clase $\mathcal{C}^1$ y, adem\'as, $\forall (x,y) \in U$ vale que
 \begin{align*}
  \dif{g}{x}(x,y) &= -\frac{\dif{F}{x} (x,y,g(x,y))}{\dif{F}{z} (x,y,g(x,y))}  \\[.4cm]
  \dif{g}{y}(x,y) &= -\frac{\dif{F}{y} (x,y,g(x,y))}{\dif{F}{z} (x,y,g(x,y))} \\
 \end{align*}
 Si $F$ es de clase $\mathcal{C}^p, p \ge 1,$ entonces $g$ tambi\'en.
\end{theorem}


\begin{theorem}\textbf{Teorema de la funci\'on impl\'icita.} \label{teo:implicit}
\mbox{}

 Sean $A \subset \Rn{n} \times \Rn{m}$ un conjunto abierto y $\mathbf{F}: A \to \Rn{m}$ una funci\'on de clase $\mathcal{C}^1$. Sea $(\mathbf{x}_0,\mathbf{z}_0) \in A$ tal que $\mathbf{F}(\mathbf{x}_0,\mathbf{z}_0) = \mathbf{0}$. Formemos el determinante
 \[
  \Delta = \det \begin{bmatrix} 
                \dif{f_1}{z_1} & \dif{f_1}{z_2} & \cdots & \dif{f_1}{z_m} \\
                               &                &        &                \\
                \dif{f_2}{z_1} & \dif{f_2}{z_2} & \cdots & \dif{f_2}{z_m} \\
                               &                &        &                \\
                \vdots         &  \vdots        & \ddots & \vdots         \\
                               &                &        &                \\
                \dif{f_m}{z_1} & \dif{f_m}{z_2} & \cdots & \dif{f_m}{z_m}
                \end{bmatrix},
 \]
 donde $\mathbf{F} = (f_1, f_2, \cdots , f_m)$ y todas las derivadas est\'an evaluadas en $(\mathbf{x}_0,\mathbf{z}_0)$. Si $\Delta \ne 0$, entonces existen un entorno abierto $U \subset \Rn{n}$ de $\mathbf{x}_0$
 %, un entorno abierto $V \subset \Rn{m}$ de $\mathbf{z}_0$ 
 y una \'unica funci\'on $\mathbf{G}: U \to \Rn{m}$ tales que $\mathbf{G}(\mathbf{x}_0) = \mathbf{z}_0$ y
 \[
  \mathbf{F}(\mathbf{x},\mathbf{G}(\mathbf{x})) = \mathbf{0},\quad \forall \mathbf{x} \in U.
 \]
 M\'as a\'un, $\mathbf{G}$ es de clase $\mathcal{C}^1$ y, adem\'as, si $\mathbf{G} = (g_1, g_2, \cdots , g_m)$,
 \[
  \begin{bmatrix} 
                \dif{g_1}{x_1} & \dif{g_1}{x_2} & \cdots & \dif{g_1}{x_n} \\
                               &                &        &                \\
                \dif{g_2}{x_1} & \dif{g_2}{x_2} & \cdots & \dif{g_2}{x_n} \\
                               &                &        &                \\
                 \vdots         &  \vdots        & \ddots & \vdots         \\
                               &                &        &                \\
                \dif{g_m}{x_1} & \dif{g_m}{x_2} & \cdots & \dif{g_m}{x_n}
                \end{bmatrix} =
 -\begin{bmatrix} 
                \dif{f_1}{z_1} & \dif{f_1}{z_2} & \cdots & \dif{f_1}{z_m} \\
                               &                &        &                \\
                \dif{f_2}{z_1} & \dif{f_2}{z_2} & \cdots & \dif{f_2}{z_m} \\
                               &                &        &                \\
                 \vdots         &  \vdots        & \ddots & \vdots         \\
                               &                &        &                \\
                \dif{f_m}{z_1} & \dif{f_m}{z_2} & \cdots & \dif{f_m}{z_m}
                \end{bmatrix}^{-1}
  \begin{bmatrix} 
                \dif{f_1}{x_1} & \dif{f_1}{x_2} & \cdots & \dif{f_1}{x_n} \\
                               &                &        &                \\
                \dif{f_2}{x_1} & \dif{f_2}{x_2} & \cdots & \dif{f_2}{x_n} \\
                               &                &        &                \\
                 \vdots         &  \vdots        & \ddots & \vdots         \\
                               &                &        &                \\
                \dif{f_m}{x_1} & \dif{f_m}{x_2} & \cdots & \dif{f_m}{x_n}
                \end{bmatrix}.            
 \]
Si $\mathbf{F}$ es de clase $\mathcal{C}^p, p \ge 1,$ entonces $\mathbf{G}$ tambi\'en.
 
\end{theorem}
